\documentclass[12pt,a4paper]{article}

% 编译建议：使用 XeLaTeX 或 LuaLaTeX，并安装完整 TeX Live（含 ctex）
% 若无 ctex，可改为 fontspec + 系统中文字体
\usepackage[UTF8]{ctex}
\usepackage{amsmath,amssymb,amsthm,amsfonts}
\usepackage{geometry}
\usepackage{array}
\usepackage{booktabs}
\usepackage{enumitem}
\usepackage{hyperref}
\usepackage{xcolor}

\geometry{margin=2.5cm}

\title{Grothendieck 构造、纤维化及相关概念}
\author{}
\date{}

\begin{document}


\maketitle

\tableofcontents
\newpage

\section*{说明}
本文档由对话内容整理而成，统一采用一个常见约定：
\begin{itemize}
  \item 基底范畴记为 \(\mathcal{B}\)；
  \item 总范畴记为 \(\mathcal{E}\)；
  \item 函子 \(p:\mathcal{E}\to\mathcal{B}\) 称为一个纤维化；
  \item 对 \(I\in\mathcal{B}\)，纤维记为
  \[
  \mathcal{E}_I=\{X\in\mathcal{E}\mid pX=I\},
  \]
  其中态射是满足 \(pf=1_I\) 的态射；
  \item “纤维化”通常指 Grothendieck fibration。
\end{itemize}

这些概念的关系可以先概括为：
\[
\boxed{
\text{indexed categories}
\;\longleftrightarrow\;
\text{Grothendieck construction}
\;\longleftrightarrow\;
\text{fibrations}
}
\]
而
\[
\boxed{
\text{cleavage}
\;\Rightarrow\;
\text{cloven fibration}
}
\]
是在一个 fibration 上选取一套“指定的笛卡尔提升”；若这套选择满足严格的复合与单位条件，就得到 split fibration。

Comprehension category 则是在 fibration 之上，再加入“上下文扩张”或“类型依赖于上下文”的额外结构。

\section{Cartesian morphism：笛卡尔态射}

设 \(p:\mathcal{E}\to\mathcal{B}\) 是一个函子。令
\[
f:X\to Y
\]
是 \(\mathcal{E}\) 中的态射，并且
\[
p(f):pX\to pY
\]
是 \(\mathcal{B}\) 中的态射。

称 \(f\) 是关于 \(p\) 的 \textbf{笛卡尔态射}，如果满足如下泛性质：

对于任意对象 \(Z\in\mathcal{E}\)、任意态射
\[
h:Z\to Y
\]
以及任意态射
\[
u:pZ\to pX
\]
只要满足
\[
p(h)=p(f)\circ u,
\]
就存在唯一的态射
\[
\bar{h}:Z\to X
\]
使得
\[
f\circ \bar{h}=h,
\qquad
p(\bar{h})=u.
\]

图示为：
\[
\begin{array}{ccc}
Z & \xrightarrow{\bar{h}} & X\\
& \searrow h & \downarrow f\\
&&Y
\end{array}
\]
并且在基底中：
\[
pZ\xrightarrow{u}pX\xrightarrow{p(f)}pY
=
pZ\xrightarrow{p(h)}pY.
\]

也就是说，\(f:X\to Y\) 对所有“从某个对象到 \(Y\) 的态射”具有一个唯一分解性质，而且这个分解在基底上必须指定为 \(u\)。

\subsection{直观理解}

设 \(p(f)=u:I\to J\)。那么 \(f\) 可以看成：
\[
f:X\to Y,
\qquad
pX=I,\quad pY=J.
\]
如果 \(f\) 是笛卡尔态射，那么 \(X\) 是 \(Y\) 沿着 \(u:I\to J\) 的“最佳拉回对象”或“逆向替换”。

因此，笛卡尔态射可以抽象地理解为：
\begin{quote}
在总范畴中实现基底态射 \(u\) 的、具有普遍性的反向运输。
\end{quote}

\subsection{纤维内态射是特殊情况}

如果 \(p(f)=1_I\)，也就是说 \(f:X\to Y\) 位于同一个纤维 \(\mathcal{E}_I\) 中，那么 \(f\) 笛卡尔的条件变成：

对任意 \(h:Z\to Y\)，如果 \(pZ=I\)，则 \(h\) 唯一地经由 \(f\) 分解。

特别地，在同一纤维中，笛卡尔态射通常只能是同构。

更准确地说：
\begin{quote}
若 \(f:X\to Y\) 是笛卡尔态射且 \(p(f)=1_{pX}\)，那么 \(f\) 是同构。
\end{quote}

因此，真正有内容的笛卡尔态射通常位于不同纤维之间。

\section{Grothendieck fibration：Grothendieck 纤维化}

函子
\[
p:\mathcal{E}\to\mathcal{B}
\]
称为一个 \textbf{Grothendieck fibration}，如果满足：

对于任意对象 \(Y\in\mathcal{E}\) 和任意基底态射
\[
u:I\to pY
\]
都存在一个笛卡尔态射
\[
\overline{u}:X\to Y
\]
使得
\[
p(\overline{u})=u.
\]

也就是说，每个对象 \(Y\) 都可以沿着任意基底态射 \(u:I\to pY\) 被“反向运输”到纤维 \(\mathcal{E}_I\)。

常把 \(\overline{u}\) 写成：
\[
u^*Y\longrightarrow Y.
\]
其中 \(u^*Y\in\mathcal{E}_I\)。

这个 \(u^*Y\) 可以理解为 \(Y\) 沿 \(u\) 的拉回、重索引或替换。

\subsection{笛卡尔提升不要求唯一选择}

Grothendieck fibration 只要求：
\begin{quote}
这样的笛卡尔提升存在。
\end{quote}
但它不要求我们为每个 \((u,Y)\) 指定一个具体的提升。

事实上，两个不同的笛卡尔提升之间存在唯一的垂直同构。也就是说，如果
\[
X\to Y
\quad\text{和}\quad
X'\to Y
\]
都是 \(u:I\to pY\) 的笛卡尔提升，那么存在唯一的同构
\[
X\cong X'
\]
位于纤维 \(\mathcal{E}_I\) 中，并且与到 \(Y\) 的两个态射相容。

因此，反向拉回 \(u^*Y\) 通常只在“唯一同构”的意义下确定，而不是严格唯一。

这正是为什么一般得到的是伪函子，而不是严格函子。

\section{Cleavage：劈分、笛卡尔提升的选择}

设 \(p:\mathcal{E}\to\mathcal{B}\) 是一个 fibration。

一个 \textbf{cleavage} 是对每一个：
\begin{itemize}
  \item 对象 \(Y\in\mathcal{E}\)；
  \item 基底态射 \(u:I\to pY\)；
\end{itemize}
选择一个笛卡尔态射
\[
\overline{u}_Y:u^*Y\to Y
\]
满足
\[
p(\overline{u}_Y)=u.
\]

换句话说，cleavage 就是：
\begin{quote}
为所有需要的笛卡尔提升，选定一个代表。
\end{quote}

通常要求恒等态射对应于恒等笛卡尔态射：
\[
(1_{pY})^*Y=Y,
\qquad
\overline{1_{pY}}_Y=1_Y.
\]
但不同文献对 cleavage 是否自动包含单位条件，约定略有不同。

\subsection{为什么需要 cleavage？}

没有 cleavage 时，对每个 \(u:I\to J\) 和 \(Y\in\mathcal{E}_J\)，只能说存在一个对象 \(u^*Y\)，并且它在同构意义下唯一。

有了 cleavage 后，才可以写出一个具体的操作：
\[
u^*:\mathcal{E}_J\to\mathcal{E}_I.
\]
对纤维内态射 \(f:Y\to Y'\)，笛卡尔泛性质给出一个唯一的态射
\[
u^*f:u^*Y\to u^*Y'.
\]
于是得到反变的“重索引函子”：
\[
u^*:\mathcal{E}_J\to\mathcal{E}_I.
\]

\section{Cloven fibration：带劈分的纤维化}

一个 \textbf{cloven fibration} 就是：
\[
p:\mathcal{E}\to\mathcal{B}
\]
再加上一个 cleavage。

所以：
\[
\boxed{
\text{cloven fibration}
=
\text{fibration}
+
\text{一个指定的 cleavage}
}
\]

必须注意：
\begin{itemize}
  \item fibration 只要求笛卡尔提升存在；
  \item cloven fibration 还指定了具体选择；
  \item 同一个 fibration 可以有很多不同的 cleavage；
  \item cleavage 通常不唯一，也不一定严格地与复合相容。
\end{itemize}

\section{Split fibration：严格可复合的劈分}

虽然用户没有单独问 split fibration，但它是理解这些概念的关键。

一个 cleavage 称为 \textbf{split cleavage}，如果满足严格的：
\[
(1_I)^*=\operatorname{id}_{\mathcal{E}_I}
\]
以及对于
\[
I\xrightarrow{u}J\xrightarrow{v}K,
\]
有
\[
(vu)^*=u^*v^*
\]
严格成立，而不只是自然同构。

因此：
\[
\boxed{
\text{split fibration}
=
\text{cloven fibration}
+
\text{严格满足单位和复合}
}
\]

一般 cleavage 只给出自然同构：
\[
u^*v^*Y\cong (vu)^*Y.
\]
所以一般 cloven fibration 对应的是伪函子，而 split fibration 对应严格函子。

\section{Grothendieck construction}

Grothendieck construction 把一个从基底范畴反变到 \(\mathbf{Cat}\) 的函子或伪函子构造成一个总范畴。

先考虑严格函子：
\[
F:\mathcal{B}^{op}\to\mathbf{Cat}.
\]
它把：
\begin{itemize}
  \item 每个 \(I\in\mathcal{B}\) 映到一个范畴 \(F(I)\)；
  \item 每个 \(u:I\to J\) 映到一个函子
  \[
  F(u):F(J)\to F(I).
  \]
\end{itemize}
注意方向是反的，所以常写作
\[
u^*=F(u):F(J)\to F(I).
\]

\subsection{Grothendieck 构造出的总范畴}

记作
\[
\int_{\mathcal{B}}F.
\]
其定义如下。

\subsubsection{对象}
对象是一个二元组
\[
(I,X),
\]
其中
\[
I\in\mathcal{B},
\qquad
X\in F(I).
\]

\subsubsection{态射}
从 \((I,X)\) 到 \((J,Y)\) 的态射由：
\begin{enumerate}
  \item 一个基底态射
  \[
  u:I\to J;
  \]
  \item 一个 \(F(I)\) 中的态射
  \[
  f:X\to F(u)(Y)
  \]
\end{enumerate}
组成。

记为
\[
(u,f):(I,X)\to(J,Y).
\]
也可以把第二部分写成：
\[
f:X\to u^*Y.
\]

\subsubsection{复合}
设
\[
(u,f):(I,X)\to(J,Y),
\]
以及
\[
(v,g):(J,Y)\to(K,Z).
\]
那么复合为
\[
(v,g)\circ(u,f)
=
\bigl(vu,\;F(u)(g)\circ f\bigr).
\]
具体地：
\[
X\xrightarrow{f}F(u)(Y)
\xrightarrow{F(u)(g)}F(u)F(v)(Z)
=
F(vu)(Z).
\]

\subsubsection{投影函子}
有一个自然的投影：
\[
\pi:\int_{\mathcal{B}}F\to\mathcal{B},
\qquad
\pi(I,X)=I,
\qquad
\pi(u,f)=u.
\]
这个投影是一个 fibration，而且由于 \(F\) 是严格函子，它自然带有一个 split cleavage。

\subsection{指定的笛卡尔态射}

给定
\[
u:I\to J
\]
和对象 \(Y\in F(J)\)，有态射
\[
(u,1_{u^*Y}):(I,u^*Y)\to(J,Y).
\]
这里
\[
u^*Y=F(u)(Y).
\]
这是一个笛卡尔态射。

所以 Grothendieck construction 中的笛卡尔提升直接由：
\[
(u,1):(I,F(u)Y)\to(J,Y)
\]
给出。

\section{Grothendieck construction 与 fibration 的对应关系}

核心定理是：
\[
\boxed{
\text{严格反变函子 }F:\mathcal{B}^{op}\to\mathbf{Cat}
\quad\Longleftrightarrow\quad
\text{split fibration }p:\mathcal{E}\to\mathcal{B}
}
\]
更准确地说，这是一个适当意义下的等价，甚至可以构成相应的 2-范畴等价。

而如果不要求严格性，则：
\[
\boxed{
\text{伪函子 }\mathcal{B}^{op}\rightsquigarrow\mathbf{Cat}
\quad\Longleftrightarrow\quad
\text{一般 fibration}
}
\]
其中：
\begin{itemize}
  \item 严格函子对应 split fibration；
  \item 伪函子对应一般 fibration；
  \item 选取 cleavage 后得到一个伪函子；
  \item 伪函子中的结合律同构对应于不同拉回次序之间的笛卡尔同构。
\end{itemize}

\subsection{从 fibration 得到纤维化伪函子}

给定 fibration
\[
p:\mathcal{E}\to\mathcal{B},
\]
每个对象 \(I\in\mathcal{B}\) 对应一个纤维范畴：
\[
I\longmapsto\mathcal{E}_I.
\]
对基底态射
\[
u:I\to J,
\]
利用 cleavage 定义：
\[
u^*:\mathcal{E}_J\to\mathcal{E}_I.
\]
但是一般来说：
\[
(1_I)^*\cong \operatorname{id}_{\mathcal{E}_I},
\]
以及
\[
u^*v^*\cong(vu)^*,
\]
这些只是自然同构，而不是严格相等。

因此得到一个伪函子：
\[
\mathcal{B}^{op}\rightsquigarrow\mathbf{Cat}.
\]

\subsection{从伪函子得到 fibration}

反过来，给定一个伪函子
\[
F:\mathcal{B}^{op}\rightsquigarrow\mathbf{Cat},
\]
也可以做类似的 Grothendieck construction：
\begin{itemize}
  \item 对象仍为 \((I,X)\)；
  \item 态射仍包含基底态射 \(u:I\to J\)；
  \item 但复合时使用伪函子的结合律同构。
\end{itemize}
所得投影
\[
\int_{\mathcal{B}}F\to\mathcal{B}
\]
是一个 fibration。

因此，Grothendieck construction 是“纤维化”和“索引范畴族”之间的桥梁。

\section{一个典型例子：切片范畴}

取范畴 \(\mathcal{C}\)，考虑对象函子
\[
\mathcal{C}^{op}\to\mathbf{Cat},
\qquad
I\longmapsto \mathcal{C}/I.
\]
对态射
\[
u:I\to J,
\]
定义拉回函子
\[
u^*:\mathcal{C}/J\to\mathcal{C}/I.
\]
如果 \(\mathcal{C}\) 有适当的拉回，这个函子把
\[
(X\to J)
\]
送到拉回：
\[
X\times_J I\to I.
\]
Grothendieck construction 的总范畴对象是：
\[
(I,X\to I).
\]
这实际上可以识别为箭头范畴
\[
\mathcal{C}^{\to}.
\]
投影
\[
\operatorname{cod}:\mathcal{C}^{\to}\to\mathcal{C}
\]
把一个箭头送到它的余域：
\[
(X\to I)\longmapsto I.
\]
这个投影是一个 fibration；笛卡尔态射正是拉回方块：
\[
\begin{array}{ccc}
X' &\to& X\\
\downarrow &&\downarrow\\
I&\xrightarrow{u}&J
\end{array}
\]
且该方块是拉回方块。

这说明“笛卡尔态射”是抽象范畴论中“拉回方块”的推广。

\section{Comprehension category：理解范畴}

Comprehension category 没有完全统一的单一定义，不同文献有不同强弱版本。最基本的思想是：
\begin{quote}
一个 fibration 描述“依赖于上下文的类型”；comprehension 结构则把一个类型变成“扩张后的上下文”。
\end{quote}

通常从一个 fibration
\[
p:\mathcal{E}\to\mathcal{B}
\]
开始。这里：
\begin{itemize}
  \item \(\mathcal{B}\) 的对象 \(\Gamma\) 被解释为上下文；
  \item 纤维 \(\mathcal{E}_\Gamma\) 的对象 \(A\) 被解释为在上下文 \(\Gamma\) 中的类型；
  \item 对象 \(A\in\mathcal{E}_\Gamma\) 应该产生一个新的上下文
  \[
  \Gamma.A;
  \]
  \item 并有一个投影
  \[
  \pi_A:\Gamma.A\to\Gamma.
  \]
\end{itemize}
这个投影通常要求是一个笛卡尔态射，且满足
\[
p(A)=\Gamma.
\]
严格来说，若总范畴中的对象 \(A\) 不是直接一个“类型对象”，通常还需要一个 comprehension 函子。

\subsection{一个常见的定义}

一种常见定义是：一个 comprehension category 是一个 fibration
\[
p:\mathcal{E}\to\mathcal{B}
\]
配备一个函子
\[
\chi:\mathcal{E}\to\mathcal{B}^{\to}
\]
使得以下三角交换：
\[
\begin{array}{ccc}
\mathcal{E} & \xrightarrow{\chi} & \mathcal{B}^{\to}\\
& \searrow p & \downarrow\operatorname{cod}\\
&&\mathcal{B}
\end{array}
\]
也就是：
\[
\operatorname{cod}(\chi(A))=p(A).
\]
如果 \(A\in\mathcal{E}_\Gamma\)，那么可写：
\[
\chi(A)=\pi_A:\Gamma.A\to\Gamma.
\]
此外通常还要求 \(\chi\) 把某类笛卡尔态射映到 \(\mathcal{B}^{\to}\) 中相应的拉回方块，或者至少保持某种与基底变化相容的结构。

不同版本可能要求：
\begin{itemize}
  \item \(\chi\) 是满射到某些箭头；
  \item \(\chi\) 在每个纤维上忠实；
  \item \(\chi\) 是纤维函子；
  \item \(\chi\) 将笛卡尔态射映成拉回方块；
  \item 只考虑满的 comprehension category；
  \item 加入终对象、单位类型、积类型、和类型、\(\Pi\)-类型等结构。
\end{itemize}
因此，看到“comprehension category”时，需要查看文献具体采用哪种定义。

\subsection{基本的 comprehension 图}

对 \(A\in\mathcal{E}_\Gamma\)，comprehension 给出：
\[
\pi_A:\Gamma.A\to\Gamma.
\]
可以把它理解为：
\begin{itemize}
  \item \(\Gamma\)：已有上下文；
  \item \(A\)：在 \(\Gamma\) 中的类型；
  \item \(\Gamma.A\)：把变量 \(x:A\) 加入上下文后的新上下文；
  \item \(\pi_A\)：忘掉新加入的变量。
\end{itemize}
在依赖类型论中，这对应于上下文扩张：
\[
\frac{\Gamma\vdash A\ \mathsf{type}}
     {\Gamma,x:A\ \mathsf{context}}.
\]

\section{Fibration 与 comprehension category 的关系}

可以这样区分：

\subsection{Fibration}
只表达：
\begin{quote}
类型可以沿上下文映射进行重索引。
\end{quote}
给定
\[
f:\Delta\to\Gamma
\]
以及
\[
A\in\mathcal{E}_\Gamma,
\]
得到
\[
f^*A\in\mathcal{E}_\Delta.
\]
这就是 substitution：
\[
A[\!f\!].
\]

\subsection{Comprehension category}
除了 substitution 之外，还表达：
\begin{quote}
每个类型 \(A\in\mathcal{E}_\Gamma\) 都可以生成一个扩张上下文 \(\Gamma.A\)。
\end{quote}
即：
\[
A\in\mathcal{E}_\Gamma
\quad\leadsto\quad
\pi_A:\Gamma.A\to\Gamma.
\]
所以：
\[
\boxed{
\text{comprehension category}
=
\text{带上下文扩张结构的 fibration}
}
\]
但注意，这只是概念性总结；严格定义要看所采用的版本。

\section{Comprehension 与 Grothendieck construction 的兼容性}

设
\[
F:\mathcal{B}^{op}\to\mathbf{Cat}
\]
是一个索引范畴族，其 Grothendieck construction 给出：
\[
p:\int_{\mathcal{B}}F\to\mathcal{B}.
\]
如果每个类型 \(A\in F(\Gamma)\) 都指定一个箭头
\[
\pi_A:\Gamma.A\to\Gamma,
\]
并且对重索引满足相容性：
\[
f^*(A)
\quad\text{对应于}\quad
\Gamma.A\times_\Gamma\Delta\to\Delta,
\]
那么这个 fibration 就具有 comprehension 结构。

也就是说，comprehension category 可以看成一种带有“对象分类/元素化”操作的 Grothendieck 构造。

\section{各概念之间的关系图}

可以把关系画成：
\[
\begin{array}{c}
F:\mathcal{B}^{op}\to\mathbf{Cat}
\\[2mm]
\Downarrow\ \text{Grothendieck construction}
\\[2mm]
\pi:\int_{\mathcal{B}}F\to\mathcal{B}
\\[2mm]
\Downarrow
\\[2mm]
\text{split fibration}
\\[2mm]
\Downarrow\ \text{忘掉严格选择}
\\[2mm]
\text{cloven fibration}
\\[2mm]
\Downarrow\ \text{忘掉 cleavage}
\\[2mm]
\text{fibration}
\end{array}
\]

另一方面：
\[
\begin{array}{c}
\text{fibration }p:\mathcal{E}\to\mathcal{B}
\\[2mm]
+\ \text{comprehension}
\\[2mm]
\Downarrow
\\[2mm]
\text{comprehension category}
\end{array}
\]

更完整地说：
\[
\begin{array}{c}
\text{strict indexed category}
\\
\Updownarrow
\\
\text{split fibration}
\end{array}
\]
\[
\begin{array}{c}
\text{pseudofunctor }\mathcal{B}^{op}\rightsquigarrow\mathbf{Cat}
\\
\Updownarrow
\\
\text{fibration}
\end{array}
\]

\section{术语间最容易混淆的地方}

\subsection{fibration 不等于 cloven fibration}
fibration：
\[
\text{笛卡尔提升存在}
\]
cloven fibration：
\[
\text{笛卡尔提升存在，并且选定了一个}
\]
所以 cleavage 是额外结构，不是 fibration 定义本身的一部分。

\subsection{cleavage 不一定满足复合律}
给定
\[
I\xrightarrow{u}J\xrightarrow{v}K,
\]
通常有两个对象：
\[
u^*(v^*X)
\]
和
\[
(vu)^*X.
\]
它们通常只满足一个规范同构：
\[
u^*v^*X\cong(vu)^*X.
\]
只有在 split cleavage 中，才要求它们严格相等。

\subsection{Grothendieck construction 的方向}
若使用
\[
F:\mathcal{B}^{op}\to\mathbf{Cat},
\]
则基底态射
\[
u:I\to J
\]
产生反向函子
\[
u^*:F(J)\to F(I).
\]
这是因为 fibration 描述的是 substitution 或 pullback：
\[
A\text{ over }J
\quad\leadsto\quad
u^*A\text{ over }I.
\]
如果使用协变函子
\[
F:\mathcal{B}\to\mathbf{Cat},
\]
相应构造通常得到 opfibration，而不是 fibration。

\subsection{Cartesian 与 cocartesian 的区别}
对于 fibration，使用的是：
\[
u:I\to J
\]
把 \(J\) 上的对象拉回到 \(I\)：
\[
u^*:\mathcal{E}_J\to\mathcal{E}_I.
\]
这是 cartesian 方向。

对于 opfibration，使用 cocartesian morphism，把对象沿着 \(u\) 正向推送：
\[
u_!:\mathcal{E}_I\to\mathcal{E}_J.
\]
因此：
\[
\text{fibration}\leftrightarrow\text{反变索引}
\]
\[
\text{opfibration}\leftrightarrow\text{协变索引}
\]

\section{一句话总结}
\begin{itemize}
  \item \textbf{Cartesian morphism}：具有泛性质的“沿基底态射反向运输”；
  \item \textbf{Grothendieck fibration}：每个基底态射都存在这样的笛卡尔提升；
  \item \textbf{Cleavage}：为每个笛卡尔提升选择一个具体代表；
  \item \textbf{Cloven fibration}：带有 cleavage 的 fibration；
  \item \textbf{Split fibration}：cleavage 还严格满足单位和复合；
  \item \textbf{Grothendieck construction}：把
  \[
  \mathcal{B}^{op}\to\mathbf{Cat}
  \]
  的索引范畴族拼成一个总范畴，其投影是 fibration；
  \item \textbf{Comprehension category}：带有“类型 \(A\) 产生上下文扩张
  \(\Gamma.A\to\Gamma\)”结构的 fibration。
\end{itemize}

最核心的对应关系是：
\[
\boxed{
\text{indexed categories}
\;\simeq\;
\text{fibrations}
}
\]
而 comprehension category 是在这个对应之上，再加入依赖类型论中的“上下文扩张”结构。

\newpage

\section{关于“拉回”的解释}

这里的“拉回”可以从两个层次理解：
\begin{enumerate}
  \item 范畴论中的 \textbf{pullback（纤维积）}；
  \item 在 fibration 中，把一个位于上下文 \(J\) 上的对象沿着 \(u:I\to J\) 重新解释为位于 \(I\) 上的对象。
\end{enumerate}

\subsection{范畴论中的拉回}

给定两个态射
\[
X\xrightarrow{f}J,
\qquad
I\xrightarrow{u}J,
\]
它们有相同的余域 \(J\)。一个拉回是对象
\[
X\times_J I
\]
连同两个投影：
\[
\begin{array}{ccc}
X\times_J I & \xrightarrow{\pi_1} & X\\
\downarrow{\pi_2} && \downarrow f\\
I & \xrightarrow{u} & J
\end{array}
\]
使得方块交换：
\[
f\circ\pi_1=u\circ\pi_2.
\]
并且满足如下泛性质：

对于任意对象 \(Z\)，如果有态射
\[
a:Z\to X,
\qquad
b:Z\to I
\]
满足
\[
f\circ a=u\circ b,
\]
那么存在唯一态射
\[
\langle a,b\rangle:Z\to X\times_J I
\]
使得
\[
\pi_1\circ\langle a,b\rangle=a,
\qquad
\pi_2\circ\langle a,b\rangle=b.
\]

图示为：
\[
\begin{array}{ccccc}
&&Z&&\\
&{}^{a}\swarrow&&\searrow^{b}&\\
X&&\longleftarrow&&I\\
&\searrow_f&&\swarrow_u&\\
&&J&
\end{array}
\]
也就是说，\(X\times_J I\) 是所有“同时映到 \(X\) 和 \(I\)，并且在 \(J\) 中相容”的对象中最普遍的一个。

\subsection{在集合中的例子}

在集合范畴 \(\mathbf{Set}\) 中，若
\[
f:X\to J,
\qquad
u:I\to J,
\]
那么拉回就是集合
\[
X\times_J I
=
\{(x,i)\in X\times I\mid f(x)=u(i)\}.
\]
所以它只保留那些在 \(J\) 中落到同一个点的配对。

例如：
\[
X=\{1,2,3\},\qquad
J=\{a,b\},
\]
定义
\[
f(1)=a,\quad f(2)=a,\quad f(3)=b.
\]
再取
\[
I=\{x,y\},
\qquad
u(x)=a,\quad u(y)=b.
\]
则
\[
X\times_J I
=
\{(1,x),(2,x),(3,y)\}.
\]
因为：
\begin{itemize}
  \item \(1,2\) 和 \(x\) 都映到 \(a\)；
  \item \(3\) 和 \(y\) 都映到 \(b\)；
  \item 其他配对在 \(J\) 中不相容。
\end{itemize}
因此，集合中的拉回就是一种“按共同索引筛选配对”的操作。

\subsection{为什么它叫“拉回”}

考虑一个对象
\[
f:X\to J.
\]
可以把它理解为“位于 \(J\) 上的对象”或“覆盖 \(J\) 的族”。

现在有一个态射
\[
u:I\to J.
\]
我们想把这个位于 \(J\) 上的对象搬到 \(I\) 上。于是构造：
\[
u^*X:=X\times_J I.
\]
这就是把 \(X\to J\) 沿着 \(u:I\to J\) 拉回：
\[
\begin{array}{ccc}
u^*X=X\times_J I &\to& X\\
\downarrow &&\downarrow\\
I&\xrightarrow{u}&J
\end{array}
\]
原来有：
\[
X\to J;
\]
拉回后得到：
\[
u^*X\to I.
\]
所以：
\[
\boxed{
u^*(X\to J)
=
(X\times_J I\to I)
}
\]
这就是“沿 \(u\) 的重索引”或“换基”。

\subsection{集合族中的拉回：最直观的解释}

设 \(X\to J\) 表示一族集合：
\[
\{X_j\}_{j\in J}.
\]
这里
\[
X_j=f^{-1}(j)
\]
是位于 \(j\) 上的那一部分。

现在给定函数
\[
u:I\to J.
\]
那么拉回后的对象
\[
u^*X=X\times_J I\to I
\]
在每个 \(i\in I\) 上的纤维是：
\[
(u^*X)_i=X_{u(i)}.
\]
也就是说，原来按 \(J\) 索引的集合族：
\[
j\longmapsto X_j
\]
经过 \(u:I\to J\) 后，变成按 \(I\) 索引的集合族：
\[
i\longmapsto X_{u(i)}.
\]
这就是“重索引”。

例如原族是：
\[
X_a=\{1,2\},\qquad X_b=\{3\}.
\]
若
\[
u:I=\{x,y,z\}\to J=\{a,b\}
\]
满足
\[
u(x)=a,\qquad u(y)=a,\qquad u(z)=b,
\]
则拉回后的族为：
\[
(u^*X)_x=X_a=\{1,2\},
\]
\[
(u^*X)_y=X_a=\{1,2\},
\]
\[
(u^*X)_z=X_b=\{3\}.
\]
所以拉回并没有改变每个元素本身，而是改变了它们所处的索引。

\subsection{在 fibration 中的“拉回”}

设
\[
p:\mathcal{E}\to\mathcal{B}
\]
是一个 fibration。

取：
\[
u:I\to J
\]
以及一个位于 \(J\) 上的对象
\[
Y\in\mathcal{E}_J.
\]
所谓沿 \(u\) 拉回 \(Y\)，就是寻找一个对象
\[
u^*Y\in\mathcal{E}_I
\]
以及一个态射
\[
\overline{u}_Y:u^*Y\to Y
\]
满足：
\[
p(\overline{u}_Y)=u,
\]
并且 \(\overline{u}_Y\) 是笛卡尔态射。

记作：
\[
\begin{array}{c}
u^*Y\\
\downarrow\\
Y
\end{array}
\]
这里要注意：
\begin{quote}
在一般 fibration 中，\(u^*Y\) 不一定是字面意义上的集合论拉回；它是通过笛卡尔泛性质定义出来的抽象“拉回”。
\end{quote}
在 \(\mathbf{Set}\) 或某些具体范畴中，它确实表现为通常的拉回；但在一般范畴中，拉回是由泛性质刻画的。

\subsection{纤维化中的笛卡尔态射为什么类似拉回方块}

考虑范畴
\[
\mathcal{C}^\to
\]
以及余域投影：
\[
\operatorname{cod}:\mathcal{C}^\to\to\mathcal{C}.
\]
一个对象是箭头：
\[
X\xrightarrow{f}J.
\]
若有基底态射
\[
u:I\to J,
\]
那么沿 \(u\) 拉回 \(f:X\to J\)，得到：
\[
\begin{array}{ccc}
X\times_J I&\to&X\\
\downarrow&&\downarrow f\\
I&\xrightarrow{u}&J
\end{array}
\]
在箭头范畴中，这对应一个态射：
\[
(X\times_J I\to I)\longrightarrow(X\to J).
\]
它在基底上的投影正是：
\[
u:I\to J.
\]
这个箭头范畴中的态射是笛卡尔态射，当且仅当原来的方块是拉回方块。

因此在这个例子中：
\[
\boxed{
\text{笛卡尔态射}
=
\text{拉回方块}
}
\]
而一般 fibration 把“拉回方块”的泛性质抽象化了。

\subsection{与依赖类型中的 substitution 的关系}

在依赖类型论中，可以把：
\begin{itemize}
  \item \(\Gamma\) 看作上下文；
  \item \(A\in\mathcal{E}_\Gamma\) 看作上下文 \(\Gamma\) 中的类型；
  \item \(f:\Delta\to\Gamma\) 看作代换。
\end{itemize}
那么拉回就是类型代换：
\[
A\in\mathcal{E}_\Gamma
\quad\leadsto\quad
f^*A\in\mathcal{E}_\Delta.
\]
通常记为：
\[
A[f].
\]
因此：
\[
\boxed{
f^*A
\text{ 就是把类型 }A\text{ 沿代换 }f\text{ 拉回到新上下文 }\Delta
}
\]
举例来说，若在上下文 \(\Gamma=(x:\mathbb{N})\) 中有类型
\[
A(x),
\]
而有代换
\[
f:\Delta\to\Gamma,
\]
例如 \(\Delta=(y:\mathbb{N})\) 且
\[
f(x)=y+1,
\]
那么拉回后的类型就是：
\[
f^*A=A(y+1).
\]
这就是 dependent type theory 中的 substitution。

\subsection{拉回的两个重要性质}

\subsubsection{恒等态射的拉回}
沿恒等态射拉回不会改变对象：
\[
1_J^*Y\cong Y.
\]
在 split fibration 中要求严格相等：
\[
1_J^*Y=Y.
\]

\subsubsection{复合的拉回}
若有：
\[
I\xrightarrow{u}J\xrightarrow{v}K,
\]
以及 \(Y\in\mathcal{E}_K\)，那么：
\[
u^*(v^*Y)\cong(vu)^*Y.
\]
意思是：
\begin{quote}
先沿 \(v\) 拉回，再沿 \(u\) 拉回，等价于直接沿复合 \(vu\) 拉回。
\end{quote}
在普通 fibration 中这是规范同构；在 split fibration 中可以要求严格相等：
\[
u^*v^*Y=(vu)^*Y.
\]

\subsection{最简洁的直观总结}

如果有一个对象“位于” \(J\) 上：
\[
Y\to J,
\]
以及一个映射：
\[
u:I\to J,
\]
那么拉回就是构造一个“位于” \(I\) 上的新对象：
\[
u^*Y\to I,
\]
使得它尽可能普遍地与原对象 \(Y\to J\) 相容。

在不同语境中，它分别表现为：
\begin{center}
\begin{tabular}{c|c}
\text{语境} & \text{拉回的含义}\\
\hline
\(\mathbf{Set}\) & \text{按函数重新索引集合族}\\
\text{几何} & \text{拉回空间、丛或截面}\\
\text{范畴论} & \text{满足泛性质的纤维积}\\
\text{fibration} & \text{笛卡尔提升}\\
\text{依赖类型论} & \text{类型沿代换进行 substitution}\\
\text{comprehension category} & \text{类型在上下文中的重索引}
\end{tabular}
\end{center}

因此在前面的讨论中，
\[
u^*Y
\]
的“拉回”本质上就是：
\begin{quote}
沿着基底态射 \(u:I\to J\)，把原本位于 \(J\) 上的对象 \(Y\) 重新解释为位于 \(I\) 上的对象。
\end{quote}

\newpage

\section{重索引函子具体是怎么得到的}

设
\[
p:\mathcal{E}\to\mathcal{B}
\]
是一个带 cleavage 的 fibration。固定一个基底态射
\[
u:I\to J.
\]
cleavage 为每个 \(Y\in\mathcal{E}_J\) 选择一个笛卡尔态射：
\[
\phi_Y^u:u^*Y\longrightarrow Y
\]
满足：
\[
p(u^*Y)=I,\qquad p(Y)=J,\qquad p(\phi_Y^u)=u.
\]
于是 \(u^*Y\) 就是 \(u^*\) 在对象 \(Y\) 上的定义。

关键是说明它如何作用在态射上。

\subsection{对象上的定义}

给定
\[
Y\in\mathcal{E}_J,
\]
cleavage 选择笛卡尔提升：
\[
\begin{array}{ccc}
u^*Y&\xrightarrow{\phi_Y^u}&Y\\
&&\downarrow p\\
I&\xrightarrow{u}&J.
\end{array}
\]
因此定义：
\[
u^*(Y):=u^*Y.
\]
现在还需要定义：
\[
u^*:\mathcal{E}_J\to\mathcal{E}_I
\]
在态射上的作用。

\subsection{对纤维内态射的定义}

取纤维 \(\mathcal{E}_J\) 中的态射：
\[
f:Y\to Y'.
\]
因为 \(f\) 位于纤维 \(\mathcal{E}_J\) 中，所以：
\[
p(f)=1_J.
\]
我们有两个已知态射：
\[
u^*Y\xrightarrow{\phi_Y^u}Y
\xrightarrow{f}Y'.
\]
其复合为：
\[
f\circ\phi_Y^u:u^*Y\to Y'.
\]
在基底中，它被 \(p\) 映成：
\[
p(f\circ\phi_Y^u)
=
p(f)\circ p(\phi_Y^u)
=
1_J\circ u
=
u.
\]
另一方面，\(\phi_{Y'}^u\) 是 \(Y'\) 沿 \(u\) 的指定笛卡尔提升：
\[
\phi_{Y'}^u:u^*Y'\to Y'.
\]
因此，根据 \(\phi_{Y'}^u\) 的笛卡尔泛性质，存在唯一态射
\[
u^*f:u^*Y\to u^*Y'
\]
使得下图交换：
\[
\begin{array}{ccc}
u^*Y & \xrightarrow{u^*f} & u^*Y'\\
\downarrow{\phi_Y^u} && \downarrow{\phi_{Y'}^u}\\
Y & \xrightarrow{f} & Y'.
\end{array}
\]
也就是：
\[
\boxed{
\phi_{Y'}^u\circ u^*f
=
f\circ\phi_Y^u
}
\]
并且由于它是纤维内态射，还满足：
\[
p(u^*f)=1_I.
\]
所以：
\[
u^*f:u^*Y\to u^*Y'
\]
确实是 \(\mathcal{E}_I\) 中的态射。

这就定义了 \(u^*\) 对态射的作用。

\subsection{图像化理解}

整个构造就是把下面这个“右边的态射”向左边提升：
\[
\begin{array}{ccccc}
u^*Y & \dashrightarrow & u^*Y'\\
\downarrow && \downarrow\\
Y&\xrightarrow{f}&Y'
\end{array}
\]
其中右边的垂直箭头
\[
u^*Y'\to Y'
\]
是笛卡尔态射，因此允许唯一地填入虚线箭头：
\[
u^*Y\to u^*Y'.
\]
所以 \(u^*f\) 不是额外随意定义的，而是由笛卡尔泛性质强制决定的。

\subsection{为什么它保持恒等态射？}

要证明：
\[
u^*(1_Y)=1_{u^*Y}.
\]
根据定义，\(u^*(1_Y)\) 是唯一满足
\[
\phi_Y^u\circ u^*(1_Y)
=
1_Y\circ\phi_Y^u
=
\phi_Y^u
\]
的态射
\[
u^*Y\to u^*Y.
\]
但恒等态射也满足：
\[
\phi_Y^u\circ 1_{u^*Y}
=
\phi_Y^u.
\]
由笛卡尔泛性质给出的唯一性，得到：
\[
\boxed{
u^*(1_Y)=1_{u^*Y}
}
\]

\subsection{为什么它保持复合？}

取纤维内态射：
\[
Y\xrightarrow{f}Y'\xrightarrow{g}Y''.
\]
我们要证明：
\[
u^*(g\circ f)=u^*g\circ u^*f.
\]
根据定义：
\[
\phi_{Y''}^u\circ u^*(g\circ f)
=
(g\circ f)\circ\phi_Y^u.
\]
另一方面：
\begin{align*}
\phi_{Y''}^u\circ(u^*g\circ u^*f)
&=
(\phi_{Y''}^u\circ u^*g)\circ u^*f\\
&=
(g\circ\phi_{Y'}^u)\circ u^*f\\
&=
g\circ(\phi_{Y'}^u\circ u^*f)\\
&=
g\circ(f\circ\phi_Y^u)\\
&=
(g\circ f)\circ\phi_Y^u.
\end{align*}
所以 \(u^*(g\circ f)\) 和 \(u^*g\circ u^*f\) 都是使上式成立的唯一态射。于是：
\[
\boxed{
u^*(g\circ f)=u^*g\circ u^*f
}
\]
因此，固定 \(u:I\to J\) 后，确实得到一个函子：
\[
\boxed{
u^*:\mathcal{E}_J\to\mathcal{E}_I
}
\]

\subsection{为什么是反变的？}

如果在基底中有：
\[
u:I\to J,
\]
那么它诱导的重索引函子方向是：
\[
u^*:\mathcal{E}_J\to\mathcal{E}_I.
\]
方向反过来了。

因此，对象层面上得到的是：
\[
I\longmapsto\mathcal{E}_I,
\]
而基底态射
\[
u:I\to J
\]
对应于：
\[
u^*:\mathcal{E}_J\to\mathcal{E}_I.
\]
所以这是一个反变的索引：
\[
\mathcal{B}^{op}\to\mathbf{Cat}.
\]

\subsection{复合时为什么通常只有自然同构？}

现在有基底态射：
\[
I\xrightarrow{u}J\xrightarrow{v}K.
\]
以及 \(Y\in\mathcal{E}_K\)。

可以有两种方式把 \(Y\) 拉回到 \(\mathcal{E}_I\)：

第一种，先沿 \(v\)，再沿 \(u\)：
\[
u^*(v^*Y).
\]
第二种，直接沿 \(vu\)：
\[
(vu)^*Y.
\]
它们通常不严格相等，但有一个规范同构：
\[
\boxed{
u^*v^*Y\cong(vu)^*Y
}
\]
为什么？因为两者到 \(Y\) 的复合都是笛卡尔态射，并且都位于基底态射 \(vu:I\to K\) 之上：
\[
u^*v^*Y\to v^*Y\to Y,
\]
以及
\[
(vu)^*Y\to Y.
\]
两个同一基底态射之上的笛卡尔提升之间存在唯一的垂直同构。

对态射也同样成立，于是得到自然同构：
\[
u^*v^*\cong(vu)^*.
\]
因此，任意带 cleavage 的 fibration 给出的通常不是严格函子
\[
\mathcal{B}^{op}\to\mathbf{Cat},
\]
而是一个伪函子：
\[
\mathcal{B}^{op}\rightsquigarrow\mathbf{Cat}.
\]

\subsection{一个具体例子：集合上的拉回}

考虑余域投影：
\[
\operatorname{cod}:\mathbf{Set}^{\to}\to\mathbf{Set}.
\]
对象是函数：
\[
q:X\to J.
\]
给定函数：
\[
u:I\to J,
\]
沿 \(u\) 的拉回对象是：
\[
u^*q:
X\times_J I\to I.
\]
现在取纤维内态射：
\[
f:(X\xrightarrow{q} J)\to(X'\xrightarrow{q'}J).
\]
它就是一个函数 \(f:X\to X'\)，满足：
\[
q'\circ f=q.
\]
那么重索引后的态射为：
\[
u^*f:X\times_J I\to X'\times_J I,
\]
定义为：
\[
u^*f(x,i)=(f(x),i).
\]
这是良定义的，因为如果
\[
q(x)=u(i),
\]
则：
\[
q'(f(x))
=
q(x)
=
u(i).
\]
所以：
\[
(f(x),i)\in X'\times_J I.
\]
这正是前面抽象构造
\[
\phi_{Y'}^u\circ u^*f
=
f\circ\phi_Y^u
\]
在集合范畴中的具体表现。

\subsection{核心结论}

给定一个 cleavage，对于每个
\[
u:I\to J
\]
和每个纤维态射
\[
f:Y\to Y',
\]
定义 \(u^*f\) 的方式是：
\[
\boxed{
u^*f
\text{ 是唯一使 }
\phi_{Y'}^u\circ u^*f
=
f\circ\phi_Y^u
\text{ 成立的态射}
}
\]
其中：
\[
\phi_Y^u:u^*Y\to Y,
\qquad
\phi_{Y'}^u:u^*Y'\to Y'
\]
是 cleavage 选定的笛卡尔提升。

笛卡尔态射的泛性质保证了：
\begin{enumerate}
  \item \(u^*f\) 存在；
  \item \(u^*f\) 唯一；
  \item \(u^*\) 保持恒等；
  \item \(u^*\) 保持复合。
\end{enumerate}
所以它自然成为一个函子：
\[
u^*:\mathcal{E}_J\to\mathcal{E}_I.
\]

\newpage

\section{“cleavage 通常不唯一，也不一定严格地与复合相容”的含义}

这句话包含两个意思：\textbf{选择不唯一}，以及\textbf{选择未必与复合严格相容}。

\subsection{“cleavage 通常不唯一”是什么意思？}

设
\[
p:\mathcal{E}\to\mathcal{B}
\]
是一个 fibration。给定
\[
u:I\to J,
\qquad
Y\in\mathcal{E}_J,
\]
fibration 只保证存在某个笛卡尔态射：
\[
\phi:u^*Y\to Y
\]
满足
\[
p(\phi)=u.
\]
但它不指定到底选哪一个。

如果还有另一个笛卡尔态射：
\[
\phi':(u^*Y)'\to Y
\]
也满足
\[
p(\phi')=u,
\]
那么 \(\phi\) 和 \(\phi'\) 都可以作为 \(Y\) 沿 \(u\) 的笛卡尔提升。

因此，一个 cleavage 就是对所有 \((u,Y)\) 选定一个提升：
\[
\phi_Y^u:u^*Y\to Y.
\]
不同的选法可能给出不同的 cleavage。

不过这些不同选择之间并不是毫无关系。两个笛卡尔提升之间有唯一的垂直同构：
\[
\theta:u^*Y\cong (u^*Y)'.
\]
这里“垂直”指：
\[
p(\theta)=1_I.
\]
所以 \(u^*Y\) 通常不是严格唯一的，而是“在唯一同构意义下唯一”。

\subsection{一个简单例子}

在集合范畴中，考虑投影：
\[
p:\mathbf{Set}^{\to}\to\mathbf{Set},
\]
即把函数
\[
X\to J
\]
送到 \(J\)。

给定：
\[
u:I\to J,
\qquad
Y=(X\to J),
\]
一个标准拉回是：
\[
X\times_J I\to I.
\]
但如果有一个集合 \(K\) 与 \(X\times_J I\) 之间的双射，那么也可以把
\[
K\to I
\]
作为同一个拉回对象的另一个表示。它们虽然不是同一个集合，但同构。

因此，严格来说，拉回对象不是唯一的；它只是在同构意义下唯一。选择哪个具体的拉回对象，就是 cleavage 中“选择”的部分。

\subsection{“不一定严格地与复合相容”是什么意思？}

设基底中有两个连续态射：
\[
I\xrightarrow{u}J\xrightarrow{v}K.
\]
设
\[
Y\in\mathcal{E}_K.
\]
我们可以有两种拉回方式。

第一种：先沿 \(v\) 拉回，再沿 \(u\) 拉回：
\[
Y
\longmapsto
v^*Y
\longmapsto
u^*(v^*Y).
\]
第二种：直接沿复合 \(v\circ u\) 拉回：
\[
Y
\longmapsto
(vu)^*Y.
\]
通常我们只有一个规范同构：
\[
\boxed{
u^*(v^*Y)\cong(vu)^*Y
}
\]
而不一定有严格相等：
\[
u^*(v^*Y)=(vu)^*Y.
\]
所谓“与复合相容”，就是指是否严格满足：
\[
\boxed{
u^*v^*=(vu)^*
}
\]
以及对恒等态射是否严格满足：
\[
\boxed{
(1_J)^*=\operatorname{id}_{\mathcal{E}_J}.
}
\]
一般 cleavage 只保证它们同构，而不保证它们字面相等。

\subsection{为什么只能得到同构？}

先沿 \(v\) 再沿 \(u\) 得到：
\[
u^*(v^*Y).
\]
它到 \(Y\) 有一个复合态射：
\[
u^*(v^*Y)
\longrightarrow
v^*Y
\longrightarrow
Y.
\]
这个复合态射位于基底态射
\[
I\xrightarrow{u}J\xrightarrow{v}K
\]
即 \(vu:I\to K\) 之上，并且是笛卡尔态射。

另一方面，cleavage 还直接选择了一个 \(vu\) 的笛卡尔提升：
\[
(vu)^*Y\to Y.
\]
于是我们有两个位于同一个基底态射 \(vu\) 之上的笛卡尔提升：
\[
u^*(v^*Y)\to Y,
\qquad
(vu)^*Y\to Y.
\]
笛卡尔提升的唯一性只告诉我们它们之间有唯一的垂直同构：
\[
u^*(v^*Y)\cong(vu)^*Y.
\]
但它并不告诉我们这两个对象必须是同一个对象。

\subsection{对应到普通拉回}

在普通范畴中，也有类似现象。

给定：
\[
X\to K,
\qquad
I\to J\to K,
\]
可以先构造：
\[
X\times_K J,
\]
再构造：
\[
(X\times_K J)\times_J I.
\]
也可以直接构造：
\[
X\times_K I.
\]
这两个对象通常不是严格相同的对象，但有自然同构：
\[
\bigl(X\times_K J\bigr)\times_J I
\cong
X\times_K I.
\]
在集合中，如果我们采用标准的有序对构造，它们甚至可能长得不同：
\begin{itemize}
  \item 一个元素可能写成 \(((x,j),i)\)；
  \item 另一个元素可能写成 \((x,i)\)。
\end{itemize}
它们显然可以自然对应，但不是字面上的同一个集合。

这正是一般 cleavage 中
\[
u^*v^*Y\cong(vu)^*Y
\]
而非严格相等的原因。

\subsection{Split cleavage 是什么？}

如果一个 cleavage 特别选择得很好，使得：
\[
(1_J)^*Y=Y
\]
并且：
\[
u^*(v^*Y)=(vu)^*Y
\]
都严格成立，那么称它是一个 \textbf{split cleavage}。

相应的 fibration 称为 \textbf{split fibration}。

所以区别是：
\begin{center}
\begin{tabular}{c|c}
\text{普通 cleavage}
&
\(u^*v^*Y\cong(vu)^*Y\)
\\[1mm]
\text{split cleavage}
&
\(u^*v^*Y=(vu)^*Y\)
\end{tabular}
\end{center}

普通 cleavage 给出的是伪函子：
\[
\mathcal{B}^{op}\rightsquigarrow\mathbf{Cat},
\]
而 split cleavage 给出严格函子：
\[
\mathcal{B}^{op}\to\mathbf{Cat}.
\]

\subsection*{总结}
“cleavage 通常不唯一”：
\begin{quote}
对同一个 \(u:I\to J\) 和 \(Y\in\mathcal{E}_J\)，可以选择不同的笛卡尔提升；它们一般只在同构意义下相同。
\end{quote}

“cleavage 不一定严格地与复合相容”：
\begin{quote}
先沿 \(v\) 再沿 \(u\) 得到的对象 \(u^*v^*Y\)，与直接沿 \(vu\) 得到的对象 \((vu)^*Y\)，通常只有规范同构，而不是严格相等。
\end{quote}

也就是：
\[
\boxed{
\text{存在性与唯一同构}
\neq
\text{严格指定且严格相等}
}
\]

\end{document}